% TOP_PROBLEM: {"id": 1800002, "title": "Eremenko–Gabrielov secant conjecture", "category_id": 6, "category": {"id": 6, "name": "geometry", "display_name": "Geometry", "description": "Euclidean and non-Euclidean geometry, geometric structures.", "slug": "geometry", "order_index": 6, "created_at": "2026-07-31T15:26:25.671Z"}, "status": "partially_solved"}
% FIRST_POSTED: null
% ENTRY_KIND: statement_only
% Imported from: batch06.tex; UnsolvedMath ID 1800002
% Compile twice: lualatex 1800002.tex
\documentclass[11pt]{article}
\usepackage{fontspec}
\setmainfont{Latin Modern Roman}
\usepackage{amsmath,amssymb,mathtools,unicode-math}
\setmathfont{Latin Modern Math}
\usepackage{xurl,hyperref}
\hypersetup{colorlinks=true,urlcolor=blue,pdftitle={UnsolvedMath Problem 1800002}}
\newcommand{\sref}[2]{\hyperlink{source:#1:#2}{[S#2]}}
\newcommand{\cataloguescope}[1]{\paragraph{Mathematical question.}#1}
\begin{document}
% UnsolvedMath ID 1800002; source catalogue code AMR-017-0002
\setcounter{section}{1800001}
\section{Eremenko–Gabrielov secant conjecture}\label{1800002}
{\small\sffamily UnsolvedMath ID 1800002 · Geometry}
\subsection{Definitions and mathematical statement}

\paragraph{Shared notation.}$\mathbb N=\{1,2,\ldots\}$, and $\mathbb Z,\mathbb Q,\mathbb R,\mathbb C$ denote the integers, rationals, reals and complex numbers. Any other conventions are specified locally.


Let $m,p\geq2$ and $d=m+p-1$. Define $F:\mathbb R\to\mathbb R^{m+p}$ by $F(x)=(1,x,\ldots,x^d)$. For $j=1,\ldots,mp$, choose pairwise disjoint real intervals $I_j$ and distinct parameters $x_{j,1},\ldots,x_{j,p}\in I_j$, and put
\[X_j=\operatorname{span}_{\mathbb C}\{F(x_{j,1}),\ldots,F(x_{j,p})\}\subset\mathbb C^{m+p}.
\]
An $m$-plane $Y\subset\mathbb C^{m+p}$ meets $X_j$ when $Y\cap X_j\neq\{0\}$, and is real when it has a basis in $\mathbb R^{m+p}$ (equivalently $\overline Y=Y$).
\paragraph{Statement recorded in the supplied source.}
Must every complex $m$-plane $Y$ meeting each $X_j$ be real? These complementary-plane incidence conditions are the divisor form of the secant conjecture, proved in the cited Theorem 1.8; no incidence conditions of higher Schubert codimension are imposed here. \sref{1800002}{2}
\subsection{Short English statement}
Are all subspaces satisfying the stated disjoint-interval secant incidences real? The cited theorem proves this particular form.
\subsection{Sources}
\begin{description}
\item[\hypertarget{source:1800002:1}{S1}] ULAM.AI and the UnsolvedMath Contributors, UnsolvedMath: A Curated Collection of Open Mathematics Problems, record 1800002 (AMR-017-0002). CC BY 4.0. \url{https://huggingface.co/datasets/ulamai/UnsolvedMath}.
\item[\hypertarget{source:1800002:2}{S2}] S. N. Karp and K. Purbhoo, \emph{Universal Plücker coordinates for the Wronski map and positivity in real Schubert calculus}, arXiv:2309.04645v2 (2026), §1.3.2, Theorem 1.8; compare §5.7.2. \url{https://arxiv.org/abs/2309.04645v2}.
\end{description}
\label{end:1800002}
\end{document}
